\documentclass[11pt,a4paper]{ctexart}
\usepackage[margin=2.2cm]{geometry}
\usepackage{xcolor}
\usepackage{enumitem}
\usepackage{amssymb}
\usepackage[normalem]{ulem}
\usepackage{titlesec}
\usepackage{hyperref}
\definecolor{navy}{HTML}{1A365D}
\definecolor{ink}{HTML}{1F2933}
\definecolor{muted}{HTML}{66788A}
\definecolor{soft}{HTML}{EEF4FA}
\definecolor{gold}{HTML}{8A5A12}
\hypersetup{colorlinks=true,linkcolor=navy,urlcolor=navy}
\setlist{nosep}
\titleformat{\section}{\Large\bfseries\color{navy}}{}{0pt}{}
\titleformat{\subsection}{\large\bfseries\color{navy}}{}{0pt}{}
\pagestyle{plain}

\title{\textbf{九下情态动词与冠词专项（网传版）}\\[4pt]
\large 人教版九年级下册英语（网传电子版）· 情态动词推测 / 冠词 · 仅供教师参考}
\author{初中英语资料站}
\date{2026-11-12}
\begin{document}
\maketitle

\section*{使用说明}
\begin{itemize}[leftmargin=1.4em]
\item ⚠ 网传版本，非官方教材，学生请勿使用，仅供教师参考。
\item 本讲义对应九下 Grammar 附录第二、三板块：情态动词推测与冠词。
\item 共 4 讲语法、6 级 48 题。
\item 在线版：\url{https://1901151458-prog.github.io/english-materials/g9x-modals/}
\end{itemize}

\begin{center}
\fbox{\parbox{0.95\textwidth}{\small\textbf{⚠ 网传版本 · 非官方教材 · 仅供教师参考}\\
本册教材尚未正式出版（官方版预计 2027 年发布），本讲义依据网传电子版整理，板块划分与题目设计可能存在偏差。\textbf{学生请勿使用}；仅供教师备课参考。正式内容请以 2027 年官方教材为准。}}
\end{center}
\bigskip
\section{must 与 can't}
\begin{center}
\begin{tabular}{p{2.4cm}p{5.4cm}p{4.4cm}}
\hline
\textbf{词} & \textbf{含义} & \textbf{例}\\
\hline
must & 一定；必定；肯定（多用于肯定句） & That movie \underline{must be} very interesting.\\
can & 可能（多用于否定句或疑问句） & \underline{Can} the news be true?\\
can't & 不可能 & That \underline{can't} be true.\\
\hline
\end{tabular}
\end{center}

\section{may / might / could}
may、might、could 都能表示「可能」。\textbf{could、might 语气更委婉，might 语气最为缓和}，表示不太确定的推测或可能性。
\begin{itemize}[leftmargin=1.4em]
\item He \underline{may / might come} to visit us tomorrow.
\item He thinks Beijing \underline{could be} the first place in Asia where humans used fire.
\end{itemize}
\textbf{推测强度：}must（一定）＞ may（可能）＞ might / could（也许，最不确定）；can't ＝ 不可能。

\section{定冠词 the}
不论是单数名词还是复数名词，也不论是可数名词还是不可数名词，其前均可用定冠词 the。主要用法：特指某人或某（些）事物；双方都知道的；上文提过的；世界上独一无二的；序数词和形容词最高级前；由普通名词构成的专有名词前；姓氏复数前；习惯用语中。
\begin{itemize}[leftmargin=1.4em]
\item \underline{The} sun is bigger than \underline{the} moon.
\item \underline{the} Atlantic Ocean / \underline{the} Wright brothers / in \underline{the} morning
\end{itemize}

\section{不定冠词与零冠词}
\begin{itemize}[leftmargin=1.4em]
\item \textbf{a / an}：a 用辅音开头前，an 用元音开头前；\textbf{只用于可数名词单数前}；用法：某一类／不具体说明／数量「一」／固定词组。
\item \textbf{零冠词}：专有名词和不可数名词前；已有 this/my/some 等定语时；复数名词表一类；星期/月份/节日前；称呼语、头衔前；一日三餐前；球类运动前；固定词组。
\end{itemize}


\newpage
\section{分级练习（48 题）}
\noindent 按难度分级排列，先做基础，再做成分与扩展，最后做综合题。

\subsection*{第 1 级 · must · 推测}
\noindent\textcolor{gray}{表示很大的可能性，「一定；必定」。}\par\smallskip
\begin{enumerate}[leftmargin=*,itemsep=4pt]
\item That movie \_\_\_ very interesting. Shall we go to see it?
\begin{enumerate}[label=\Alph*.,leftmargin=2em,itemsep=1pt,topsep=2pt]
\item must be
\item can be
\item can’t be
\item might be
\end{enumerate}
\item The cakes look nice. They \_\_\_ delicious.
\begin{enumerate}[label=\Alph*.,leftmargin=2em,itemsep=1pt,topsep=2pt]
\item must be
\item can be
\item can’t be
\item need be
\end{enumerate}
\item must 表推测时意为 \_\_\_。
\begin{enumerate}[label=\Alph*.,leftmargin=2em,itemsep=1pt,topsep=2pt]
\item 一定；必定；肯定
\item 可能
\item 不可能
\item 必须
\end{enumerate}
\item must 表推测多用于 \_\_\_。
\begin{enumerate}[label=\Alph*.,leftmargin=2em,itemsep=1pt,topsep=2pt]
\item 肯定句
\item 否定句
\item 疑问句
\item 条件句
\end{enumerate}
\item He looks tired. He \_\_\_ have worked all night.
\begin{enumerate}[label=\Alph*.,leftmargin=2em,itemsep=1pt,topsep=2pt]
\item must
\item can
\item can’t
\item might not
\end{enumerate}
\item The light is on. Someone \_\_\_ be at home.
\begin{enumerate}[label=\Alph*.,leftmargin=2em,itemsep=1pt,topsep=2pt]
\item must
\item can’t
\item may not
\item needn’t
\end{enumerate}
\item 下面哪个词表示「一定」？
\begin{enumerate}[label=\Alph*.,leftmargin=2em,itemsep=1pt,topsep=2pt]
\item must
\item can
\item can’t
\item might
\end{enumerate}
\item It \_\_\_ be delicious. It smells so good.
\begin{enumerate}[label=\Alph*.,leftmargin=2em,itemsep=1pt,topsep=2pt]
\item must
\item can’t
\item may not
\item needn’t
\end{enumerate}
\end{enumerate}

\subsection*{第 2 级 · can / can’t · 推测}
\noindent\textcolor{gray}{can 表可能（否/疑），can’t 表不可能。}\par\smallskip
\begin{enumerate}[leftmargin=*,itemsep=4pt]
\item — \_\_\_ the news be true?
\begin{enumerate}[label=\Alph*.,leftmargin=2em,itemsep=1pt,topsep=2pt]
\item Can
\item Must
\item May
\item Need
\end{enumerate}
\item That \_\_\_ be true. I don’t believe it.
\begin{enumerate}[label=\Alph*.,leftmargin=2em,itemsep=1pt,topsep=2pt]
\item can’t
\item must
\item may
\item should
\end{enumerate}
\item It looks so thin. It \_\_\_ be heavy.
\begin{enumerate}[label=\Alph*.,leftmargin=2em,itemsep=1pt,topsep=2pt]
\item can’t
\item must
\item may
\item should
\end{enumerate}
\item can 表「可能」常用于 \_\_\_。
\begin{enumerate}[label=\Alph*.,leftmargin=2em,itemsep=1pt,topsep=2pt]
\item 否定句或疑问句
\item 肯定句
\item 祈使句
\item 感叹句
\end{enumerate}
\item can’t 表示 \_\_\_。
\begin{enumerate}[label=\Alph*.,leftmargin=2em,itemsep=1pt,topsep=2pt]
\item 不可能
\item 一定
\item 可能
\item 必须
\end{enumerate}
\item can 用于疑问句表推测时，常带有 \_\_\_ 的含义。
\begin{enumerate}[label=\Alph*.,leftmargin=2em,itemsep=1pt,topsep=2pt]
\item 惊异；怀疑；不相信
\item 肯定
\item 命令
\item 礼貌
\end{enumerate}
\item He \_\_\_ be at home. He has gone to Beijing.
\begin{enumerate}[label=\Alph*.,leftmargin=2em,itemsep=1pt,topsep=2pt]
\item can’t
\item must
\item may
\item should
\end{enumerate}
\item \_\_\_ it be true that he won?
\begin{enumerate}[label=\Alph*.,leftmargin=2em,itemsep=1pt,topsep=2pt]
\item Can
\item Must
\item Should
\item Need
\end{enumerate}
\end{enumerate}

\subsection*{第 3 级 · may / might / could}
\noindent\textcolor{gray}{不太确定的推测，might 最缓和。}\par\smallskip
\begin{enumerate}[leftmargin=*,itemsep=4pt]
\item He \_\_\_ come to visit us tomorrow.
\begin{enumerate}[label=\Alph*.,leftmargin=2em,itemsep=1pt,topsep=2pt]
\item may
\item must
\item can’t
\item need
\end{enumerate}
\item He thinks Beijing \_\_\_ be the first place in Asia where humans used fire.
\begin{enumerate}[label=\Alph*.,leftmargin=2em,itemsep=1pt,topsep=2pt]
\item could
\item must
\item can’t
\item need
\end{enumerate}
\item The evidence shows that they \_\_\_ also have developed the skill of making silk.
\begin{enumerate}[label=\Alph*.,leftmargin=2em,itemsep=1pt,topsep=2pt]
\item might
\item must
\item can’t
\item need
\end{enumerate}
\item may、might、could 中，语气最缓和的是 \_\_\_。
\begin{enumerate}[label=\Alph*.,leftmargin=2em,itemsep=1pt,topsep=2pt]
\item might
\item may
\item could
\item 都一样
\end{enumerate}
\item could、might 表示推测时比 may 语气 \_\_\_。
\begin{enumerate}[label=\Alph*.,leftmargin=2em,itemsep=1pt,topsep=2pt]
\item 更委婉
\item 更强烈
\item 更肯定
\item 完全否定
\end{enumerate}
\item It \_\_\_ rain tomorrow, but I’m not sure.
\begin{enumerate}[label=\Alph*.,leftmargin=2em,itemsep=1pt,topsep=2pt]
\item might
\item must
\item can’t
\item needn’t
\end{enumerate}
\item 下面哪个词表示「不太确定的推测」？
\begin{enumerate}[label=\Alph*.,leftmargin=2em,itemsep=1pt,topsep=2pt]
\item might
\item must
\item can’t
\item need
\end{enumerate}
\item She \_\_\_ be in the library. I saw her there just now.
\begin{enumerate}[label=\Alph*.,leftmargin=2em,itemsep=1pt,topsep=2pt]
\item could
\item can’t
\item must not
\item needn’t
\end{enumerate}
\end{enumerate}

\subsection*{第 4 级 · 定冠词 the}
\noindent\textcolor{gray}{特指、独一无二、序数词与最高级等。}\par\smallskip
\begin{enumerate}[leftmargin=*,itemsep=4pt]
\item \_\_\_ sun is bigger than \_\_\_ moon.
\begin{enumerate}[label=\Alph*.,leftmargin=2em,itemsep=1pt,topsep=2pt]
\item The; the
\item A; a
\item 不填; 不填
\item The; a
\end{enumerate}
\item Da Vinci was one of \_\_\_ first Western artists to study the human body.
\begin{enumerate}[label=\Alph*.,leftmargin=2em,itemsep=1pt,topsep=2pt]
\item the
\item a
\item an
\item 不填
\end{enumerate}
\item The blue whale is \_\_\_ largest and the heaviest animal.
\begin{enumerate}[label=\Alph*.,leftmargin=2em,itemsep=1pt,topsep=2pt]
\item the
\item a
\item an
\item 不填
\end{enumerate}
\item \_\_\_ Wright brothers are famous.
\begin{enumerate}[label=\Alph*.,leftmargin=2em,itemsep=1pt,topsep=2pt]
\item The
\item A
\item An
\item 不填
\end{enumerate}
\item \_\_\_ Atlantic Ocean is very large.
\begin{enumerate}[label=\Alph*.,leftmargin=2em,itemsep=1pt,topsep=2pt]
\item The
\item A
\item An
\item 不填
\end{enumerate}
\item Let’s go and buy something for \_\_\_ party.
\begin{enumerate}[label=\Alph*.,leftmargin=2em,itemsep=1pt,topsep=2pt]
\item the
\item a
\item an
\item 不填
\end{enumerate}
\item \_\_\_ file is sent to the 3D printer.
\begin{enumerate}[label=\Alph*.,leftmargin=2em,itemsep=1pt,topsep=2pt]
\item The
\item A
\item An
\item 不填
\end{enumerate}
\item We usually meet in \_\_\_ morning.
\begin{enumerate}[label=\Alph*.,leftmargin=2em,itemsep=1pt,topsep=2pt]
\item the
\item a
\item an
\item 不填
\end{enumerate}
\end{enumerate}

\subsection*{第 5 级 · 不定冠词 a / an}
\noindent\textcolor{gray}{某一种类、不具体说明、数量「一」。}\par\smallskip
\begin{enumerate}[leftmargin=*,itemsep=4pt]
\item I want to be \_\_\_ engineer in the future.
\begin{enumerate}[label=\Alph*.,leftmargin=2em,itemsep=1pt,topsep=2pt]
\item an
\item a
\item the
\item 不填
\end{enumerate}
\item It’s \_\_\_ big farm.
\begin{enumerate}[label=\Alph*.,leftmargin=2em,itemsep=1pt,topsep=2pt]
\item a
\item an
\item the
\item 不填
\end{enumerate}
\item \_\_\_ true friend is always there for you.
\begin{enumerate}[label=\Alph*.,leftmargin=2em,itemsep=1pt,topsep=2pt]
\item A
\item An
\item The
\item 不填
\end{enumerate}
\item Our town has just built \_\_\_ new hospital and library.
\begin{enumerate}[label=\Alph*.,leftmargin=2em,itemsep=1pt,topsep=2pt]
\item a
\item an
\item the
\item 不填
\end{enumerate}
\item 不定冠词只能用于 \_\_\_ 前。
\begin{enumerate}[label=\Alph*.,leftmargin=2em,itemsep=1pt,topsep=2pt]
\item 可数名词单数
\item 可数名词复数
\item 不可数名词
\item 任何名词
\end{enumerate}
\item \_\_\_ hour 是以元音音素开头，用 \_\_\_。
\begin{enumerate}[label=\Alph*.,leftmargin=2em,itemsep=1pt,topsep=2pt]
\item an
\item a
\item the
\item 不填
\end{enumerate}
\item \_\_\_ English teacher（一位英语老师）
\begin{enumerate}[label=\Alph*.,leftmargin=2em,itemsep=1pt,topsep=2pt]
\item an
\item a
\item the
\item 不填
\end{enumerate}
\item There is \_\_\_ moment, please wait.
\begin{enumerate}[label=\Alph*.,leftmargin=2em,itemsep=1pt,topsep=2pt]
\item a
\item an
\item the
\item 不填
\end{enumerate}
\end{enumerate}

\subsection*{第 6 级 · 零冠词与综合}
\noindent\textcolor{gray}{不用冠词的场合与三类冠词混辨。}\par\smallskip
\begin{enumerate}[leftmargin=*,itemsep=4pt]
\item \_\_\_ China is a great country.
\begin{enumerate}[label=\Alph*.,leftmargin=2em,itemsep=1pt,topsep=2pt]
\item 不填
\item A
\item An
\item The
\end{enumerate}
\item I go to school \_\_\_ Monday.
\begin{enumerate}[label=\Alph*.,leftmargin=2em,itemsep=1pt,topsep=2pt]
\item on
\item in
\item at
\item the
\end{enumerate}
\item For \_\_\_ lunch, she likes hamburgers and salad.
\begin{enumerate}[label=\Alph*.,leftmargin=2em,itemsep=1pt,topsep=2pt]
\item 不填
\item a
\item an
\item the
\end{enumerate}
\item They often play \_\_\_ football after school.
\begin{enumerate}[label=\Alph*.,leftmargin=2em,itemsep=1pt,topsep=2pt]
\item 不填
\item a
\item an
\item the
\end{enumerate}
\item These ping-pong bats are nice. 此句中 these 的作用相当于 \_\_\_。
\begin{enumerate}[label=\Alph*.,leftmargin=2em,itemsep=1pt,topsep=2pt]
\item 零冠词（前已有定语）
\item 定冠词
\item 不定冠词
\item 冠词不可替代
\end{enumerate}
\item He was \_\_\_ very tired. He \_\_\_ have worked all day.
\begin{enumerate}[label=\Alph*.,leftmargin=2em,itemsep=1pt,topsep=2pt]
\item 不填; must
\item the; can
\item a; can’t
\item an; might
\end{enumerate}
\item This is \_\_\_ Professor Liu. 此空应填 \_\_\_。
\begin{enumerate}[label=\Alph*.,leftmargin=2em,itemsep=1pt,topsep=2pt]
\item 不填
\item a
\item an
\item the
\end{enumerate}
\item 那栋楼不可能是新的。The building \_\_\_ be new.
\begin{enumerate}[label=\Alph*.,leftmargin=2em,itemsep=1pt,topsep=2pt]
\item can’t
\item must
\item may
\item should
\end{enumerate}
\end{enumerate}

\newpage
\section*{答案与解析}
\small
\begin{enumerate}[leftmargin=1.6em,itemsep=2pt]
\item \textbf{M1-1 A}：must 表示很大可能性「一定」；教材例句。
\item \textbf{M1-2 A}：must 表肯定推测；教材例句。
\item \textbf{M1-3 A}：教材：must 表示很大的可能性，意为「一定；必定；肯定」。
\item \textbf{M1-4 A}：教材：多用于肯定句中。
\item \textbf{M1-5 A}：很有把握的推测用 must。
\item \textbf{M1-6 A}：有把握的肯定推测用 must。
\item \textbf{M1-7 A}：must 表示一定。
\item \textbf{M1-8 A}：根据迹象做出的肯定推测用 must。
\item \textbf{M2-1 A}：can 用于疑问句表可能，带惊异/怀疑；教材例句。
\item \textbf{M2-2 A}：can’t 表示不可能；教材例句。
\item \textbf{M2-3 A}：can’t 表示不可能；教材例句。
\item \textbf{M2-4 A}：教材：can 表示「可能」，常用于否定句或疑问句中。
\item \textbf{M2-5 A}：教材：can’t 表示「不可能」。
\item \textbf{M2-6 A}：教材原话。
\item \textbf{M2-7 A}：已去北京，所以不可能在家，用 can’t。
\item \textbf{M2-8 A}：疑问句表可能用 can。
\item \textbf{M3-1 A}：may 表示不太确定的推测；教材例句。
\item \textbf{M3-2 A}：could 表示委婉的不确定推测；教材例句。
\item \textbf{M3-3 A}：might 表示推测；教材例句。
\item \textbf{M3-4 A}：教材：might 语气最为缓和。
\item \textbf{M3-5 A}：教材：could、might 语气更委婉。
\item \textbf{M3-6 A}：不确定的推测用 might。
\item \textbf{M3-7 A}：might 表示不太确定的推测或可能性。
\item \textbf{M3-8 A}：could 表示不确定的推测（可能）。
\item \textbf{M4-1 A}：世界上独一无二的事物前用 the；教材表2例句。
\item \textbf{M4-2 A}：序数词前用 the；教材表2例句。
\item \textbf{M4-3 A}：形容词最高级前用 the；教材表2例句。
\item \textbf{M4-4 A}：姓氏的复数形式前用 the；教材表2。
\item \textbf{M4-5 A}：由普通名词构成的专有名词前用 the；教材表2例句。
\item \textbf{M4-6 A}：特指某个 party，用 the；教材表2例句。
\item \textbf{M4-7 A}：指上文提过的事物，用 the；教材表2例句。
\item \textbf{M4-8 A}：in the morning 是习惯用语；教材表2。
\item \textbf{M5-1 A}：engineer 以元音音素开头，用 an；教材表3例句。
\item \textbf{M5-2 A}：big 以辅音音素开头，用 a；教材表3例句。
\item \textbf{M5-3 A}：指某一类人中的一个，用 a；教材表3例句。
\item \textbf{M5-4 A}：表示数量「一」；教材表3例句。
\item \textbf{M5-5 A}：教材：不定冠词只用于可数名词单数前。
\item \textbf{M5-6 A}：hour 的 h 不发音，以元音音素开头用 an；教材表3例句。
\item \textbf{M5-7 A}：English 以元音音素开头用 an；教材表3例句。
\item \textbf{M5-8 A}：a moment 属固定词组；教材表3。
\item \textbf{M6-1 A}：专有名词前用零冠词；教材表4。
\item \textbf{M6-2 A}：星期前用零冠词 + 介词 on；教材表4例句 Monday。
\item \textbf{M6-3 A}：一日三餐前用零冠词；教材表4例句。
\item \textbf{M6-4 A}：球类运动名称前用零冠词；教材表4。
\item \textbf{M6-5 A}：名词前已有 this/these 等定语时用零冠词；教材表4。
\item \textbf{M6-6 A}：tired 前零冠词；有把握推测用 must。
\item \textbf{M6-7 A}：称呼语、头衔和职务名称前用零冠词；教材表4例句。
\item \textbf{M6-8 A}：不可能用 can’t。
\end{enumerate}

\end{document}
